\section{Discussion}
\label{sec:discussion}


The initial campaign establishes \SHARP\ as a
wide-area source of repeated main-belt asteroid
measurements. The population represented by the
primary associations, together with the observed
elongation distribution, locates the survey's
current scientific emphasis. Its most immediate
contribution is additional astrometric and
photometric sampling of cataloged asteroids
over a broad ecliptic region.

The footprint and cadence characterize
complementary aspects of that contribution.
Nearly all of the accumulated geometric
footprint received three exposures within a
single night at some point in the campaign,
whereas only 67.43\% of individual
field--night pointing groups supplied at
least three validated visits. Reaching a
field on one well-sampled night does not
ensure a complete sequence whenever the
field is revisited. Completion of repeated
visits is therefore an important factor in
the area available for constructing
trajectories on a particular night.

At the median adjacent-visit interval of
37.65~min, the historical search window
corresponds to a displacement of approximately
$1.9$--$39.5\asec$. At the slow end,
adjacent-visit displacements are only
a few arcseconds, making the treatment
of stationary-source matches and
small-displacement pairs particularly
relevant. The reduced conditional retention
at $3$--$5$~\rateunit\ identifies a part of
the detected input sample for which the
stored processing is less effective, although
the present aggregate test does not isolate
the responsible processing stage.


The known-object sample combines more than
half a million positional associations with
calibrated aperture photometry. Its
object-level revisit distributions provide
the information needed to select targets
with useful time coverage. The median
four measurements on two nights describes
sparse sampling for a typical identity,
while the more frequently observed objects
offer longer sequences for subsequent
analysis. Rotational studies will require
object-by-object assessment of phase
coverage, blending, changing observing
geometry, and the photometric passband
\citep{WarnerHarrisPravec2009}.

The measured photometric repeatability
and astrometric residuals support
using these catalogs to construct
selected time series with documented
uncertainties. Cadence remains an
equally important constraint: completing
three visits supplies a within-night
trajectory, whereas observations spread
across additional nights extend the
astrometric arc and the range of
rotational phases sampled.

Population-level applications will
also require a selection function.
The conditional retention test
describes linking after catalog
detection. A total selection function
would additionally need the
probability of obtaining a catalog
detection and a usable sequence
for each source and observing
condition.


The retained candidate sample supplies
positions, angular motion, and brightness
over observed single-night arcs. The
short-arc fits establish a consistency
check, but longer baselines are needed
to constrain the poorly determined
line-of-sight motion and obtain secure
identities \citep{Kubica2007,Holman2018,Moeyens2021}.
Connecting existing detections across
nights and obtaining timely follow-up
would therefore add physical information
to the single-night sample.

The C/2025~Y1 example also shows how the
contents of the reference orbit catalog
affect the interpretation of a residual
moving-source stream. Maintaining
asteroid and comet attribution as
explicit steps will allow future
searches to distinguish catalog coverage
from unexplained motion. The initial
results provide concrete trajectories
on which those developments can be
evaluated.
